\title{Large Language Models Can Automatically Engineer Features for Few-Shot Tabular Learning}

\begin{abstract}
Large Language Models (LLMs) have exhibited extraordinary proficiency in reasoning through complex, previously unseen challenges, making them highly promising for tabular data learning—a cornerstone for many practical deployments. This paper introduces a new in-context learning paradigm, \textsf{FeatLLM}, wherein LLMs are tasked with engineering features, thereby transforming raw input data into representations that are particularly advantageous for predictive modeling with tabular data. These features facilitate accurate class probability estimation when used with simple downstream models such as linear regression, yielding strong results in few-shot scenarios. Importantly, \textsf{FeatLLM} confines its reliance on such lightweight predictive models, utilizing only the pre-computed features at inference, and foregoes repeated LLM querying per inference sample, in contrast to prior LLM-based strategies. This setup only necessitates API-level LLM access, circumventing prompt size restrictions. Through extensive empirical validation on a variety of tabular benchmarks, we show that \textsf{FeatLLM} generates superior rules, achieving substantial improvements—on average 10\%—over established methods like TabLLM and STUNT.
\end{abstract}

\section{Introduction}
Large Language Models (LLMs) have, in recent years, demonstrated remarkable competence in generating coherent responses for previously unencountered tasks, without explicit parameter adjustment for each task~\cite{creswell2022selection,imani2023mathprompter,taylor2022galactica}. Their exposure to massive textual datasets endows LLMs with extensive implicit knowledge, allowing them to substitute for domain experts across numerous contexts~\cite{Genomes2021,singhal2023large,wilcox2003role}. As a result, LLMs are central to the automation of tasks spanning dialogue understanding~\cite{finch2023leveraging} to program repair~\cite{fan2023automated}. Significantly, by furnishing task descriptions and a limited number of examples as part of the prompt, LLMs interpret the task context and appropriately prioritize relevant variables and samples~\cite{brown2020language,zhao2023survey}. This context-driven analysis signals the potential for LLMs to serve as proficient feature engineers.

The exploitation of LLMs’ embedded knowledge and logical reasoning for tabular data analysis has recently gained traction. For example, leveraging knowledge about the heightened risk of certain medical conditions among older adults can meaningfully augment disease prediction systems. In this setting, recent approaches articulate both task and feature information in natural language, serialize inputs accordingly, and utilize the resulting representations as input to LLMs~\cite{dinh2022lift,wang2023anypredict}. Subsequently, either in-context learning~\cite{nam2023semi} or efficient fine-tuning methods~\cite{dinh2022lift,hegselmann2023tabllm} are applied. The capacity of LLMs to incorporate prior knowledge has been shown to substantially surpass standard tabular learning methods, particularly in situations where data are limited~\cite{hegselmann2023tabllm}.

Existing tabular learning methods that utilize LLMs face several notable drawbacks. To perform end-to-end prediction tasks, they require at least one inference per input, leading to considerable computational demands. Additionally, many techniques depend on fine-tuning LLMs to achieve competitive accuracy, which becomes problematic when the underlying LLM does not provide a public training pipeline or tuning API—an increasingly common limitation as state-of-the-art LLMs restrict access solely to inference APIs~\cite{anil2023palm,openai2023gpt4}. Moreover, most approaches are ill-equipped to handle extended prompt lengths~\cite{liu2023lost}; as tabular data grows in feature count and text serialization increases, their applicability may diminish or performance may degrade. These issues fundamentally hinder the practical deployment of LLM-driven tabular learning solutions.

In contrast, our method departs from the typical reliance on LLMs for direct predictions. Instead, we repurpose LLMs as feature engineers, harnessing their prior knowledge to generate informative features efficiently. Concretely, we introduce a new in-context learning strategy, termed \textsf{FeatLLM}, which focuses on extracting the ``criteria'' governing LLM decisions. For example, in a disease prediction scenario, the LLM can articulate explicit rules that specify which feature combinations lead to positive identification. These extracted criteria or rules are leveraged to construct new, more effective features, which then serve as replacements for original features in downstream tasks. This reframing delegates feature construction to the LLM, simplifying subsequent modeling—once salient features are generated, a straightforward model suffices for inference, thus reducing latency. Utilizing the LLMs’ in-context learning capacity, features may be distilled directly by supplying examples in the input prompt, obviating any training requirement. To address prompt-length constraints, we further employ ensemble techniques such as feature bagging~\cite{breiman2001random}, enabling more robust operation despite large or complex feature sets.

\textsf{FeatLLM} organizes the prompt for feature engineering into two distinct phases: (1) comprehending the given problem and deducing the interplay between features and targets; and (2) leveraging the insights gathered as well as few-shot sample data to formulate clear, class-separating rules. This methodical reasoning process encourages LLMs to disregard irrelevant features during rule formulation. The derived rules are subsequently utilized on the dataset (interpreted as code), converting the data into binary indicators reflecting the compliance of each sample with the individual rules. A simple model is then trained to predict class probabilities using these newly generated binary features. To enhance reliability, the approach incorporates ensembles by producing features repeatedly in combination with feature bagging. By tuning the counts of features and samples selected during bagging, the issue of overly large prompts is alleviated. \textsf{FeatLLM} does not require any training phase and imposes minimal inference costs, making LLM-based applications practical across diverse real-world contexts.

We assess \textsf{FeatLLM} across 13 tabular datasets under limited-data conditions, demonstrating its consistent and high performance. The experiments indicate that LLMs effectively draw on both prior and empirical knowledge to extract meaningful features. Our approach surpasses state-of-the-art few-shot learning methods in multiple scenarios. The implementation is available via an anonymized GitHub repository at~\texttt{https://github.com/Sungwon-Han/FeatLLM}.

\section{Related Work}

\subsection{Few-Shot Learning with Tabular Data} Few-shot learning has revolutionized the way general representations are derived from datasets, significantly reducing the dependency on extensive labeling~\cite{chen2018closer}. While the field’s origins are rooted in image data, few-shot paradigms have exhibited robust adaptability across domains, including text, audio, and lately, tabular data~\cite{majumder2022few,nam2023stunt,schick2021exploiting}. The challenge of learning with limited labeled data is the propensity for models to pick up on irrelevant correlations~\cite{bartlett2021deep}. Thus, contemporary approaches integrate supplementary information or apply domain-driven prior knowledge, embedding suitable inductive biases during training. Examples include leveraging abundant unlabeled data with careful data augmentation strategies within semi-supervised frameworks~\cite{ucar2021subtab,yoon2020vime}, or achieving effective model initialization through unsupervised, task-independent pre-training~\cite{somepalli2021saint}. Unsupervised methods can involve masking or introducing noise to sections of the input, followed by optimizing reconstruction-based objectives~\cite{arik2021tabnet,majmundar2022met}, or adopting data augmentation for contrastive learning~\cite{bahri2022scarf,chen2020simple}. The diversity inherent in tabular data, however, complicates the search for universally effective augmentations, limiting the impact of these techniques in few-shot contexts~\cite{nam2023stunt}. 

More recently, the focus has shifted towards training models across an extensive array of tabular datasets—transcending direct relevance to the target task—and transferring the acquired knowledge via transfer learning~\cite{zhang2023generative,levin2022transfer,zhu2023xtab}. TransTab~\cite{wang2022transtab}, for example, utilizes transformer architectures to model semantic relationships among columns through their names, beyond mere tabular values. Insights gleaned from multiple tables and column semantics have enabled zero-shot and few-shot generalization to unseen tabular problems. Another approach, TabPFN~\cite{hollmann2022tabpfn}, synthesizes varied distributions to generate artificial data for Bayesian neural network pre-training. Subsequent inference involves directly inputting both training and test instances from the target task, bypassing further training phases. The broad-based prior learned from diverse datasets allows these models to surpass classical tree-based algorithms and proves especially advantageous for few-shot scenarios.

\subsection{Language-Interfaced Tabular Learning} Integrating large language models (LLMs)~\cite{anil2023palm,openai2023gpt4} represents a further avenue for leveraging prior knowledge. LLMs, thanks to their exposure to vast and heterogeneous textual datasets, are capable of extrapolating to unfamiliar tasks, thereby showing cross-domain versatility~\cite{brown2020language}. Contemporary efforts utilize the extensive prior embedded within LLMs to enhance tabular learning, often by converting tabular structures into text and using these textual representations as prompts~\cite{dinh2022lift,hegselmann2023tabllm,wang2023anypredict}. Including tables, feature metadata, and task context in prompts enables LLMs to deduce linkages among features and tasks for inference purposes. Advances span from training models with methods aimed at efficient parameter adaptation, such as LoRA~\cite{hu2021lora} or IA3~\cite{liu2022few}, to employing in-context learning, where a handful of example tasks are embedded in prompts without explicit model training~\cite{dinh2022lift,hegselmann2023tabllm,nam2023semi,slack2023tablet}.

Although previous techniques have proven successful in enhancing few-shot tabular learning, their reliance on consistently routing inference through the LLM limits their practicality in industrial settings. The approach taken in this work intentionally moves away from standard end-to-end, black-box sample processing with LLMs. Rather than predicting each sample via the LLM, the model focuses the LLM’s capabilities on discerning rules or conditions linked to each class. \textsf{FeatLLM} interprets these identified rules as programmatic code to generate new feature representations. This enables downstream predictions to occur independently from the LLM, exploiting in-context learning to derive rules using prior knowledge and few-shot sample information.

\section{Methods}
\textsf{FeatLLM} is developed to retrieve “rules”—class-conditional criteria—from its limited sample inputs, as illustrated in Figure~\ref{fig:main_model}. These conditions are extracted by the LLM, which then parses the outputs and synthesizes new binary features for each input sample, signaling whether the corresponding rule holds. These features are subsequently aggregated using a dot product with a set of trainable, non-negative weights to compute class probabilities. Through model training, the algorithm uncovers the relative significance of each extracted feature in relation to class prediction (see Section~\ref{Sec:training} for details). Multiple rounds of this procedure are conducted in a bagging framework, and the results are amalgamated using ensemble strategies. The problem setup and the particulars of each step are elaborated as follows.

\vspace*{-0.12in}
\paragraph{Problem formulation.} We begin with a tabular dataset containing $N$ annotated samples, denoted as $\mathcal{D} = \{(\mathbf{x}^i, \mathbf{y}^i)\}_{i=1}^{N}$. Each sample within $\mathcal{D}$ comprises $d$ distinct features—here, each $\mathbf{x}^i$ represents a vector of $d$ dimensions—and these features are accompanied by their respective names in natural language, captured as $F = \{f_j\}_{j=1}^{d}$, as well as definitions stored separately. In this context, assuming a classification scenario, each label $\mathbf{y}^i$ is drawn from a finite set of categories. For $k$-shot learning settings, only $k$ samples (where $k<N$) are randomly selected and used for training purposes.

\subsection{Prompt Design for Extracting Rules}
\label{Sec:prompt_design}
To facilitate the extraction of rules by the LLM via a reasoning process that closely parallels expert human problem-solving strategies in tabular learning, we construct prompts consisting of three core elements (see Fig.~\ref{fig:prompt_for_rule} and Appendix~\ref{sec:example_rule}):

\vspace*{-0.12in}
\paragraph{Basic information description.} This section conveys fundamental knowledge required for the task at hand (illustrated by the orange segments in Fig.~\ref{fig:prompt_for_rule}). It features a clear statement of the task (complete with label details), descriptions of input features, as well as sample demonstrations. The task is articulated as a question format (such as ``Does this patient's myocardial infarction complication data indicate chronic heart failure? Yes or no?"), with further examples shown in Appendix~\ref{sec:task_desc}. Feature descriptions specify value types and optionally clarify their definitions; for categorical features, allowable categories may be enumerated as examples. To illustrate, several labeled training samples are transformed into textual form and presented together with their associated ground-truth outputs. This serialization process mirrors the conventions adopted in previous works~\cite{dinh2022lift,hegselmann2023tabllm}:

\begin{align}
&\text{Serialize}(\mathbf{x}^i,\mathbf{y}^i, F) = \nonumber \\ 
&``\ f_1\ \text{is}\ \mathbf{x}^i_1.\ ... \ f_d\  \text{is}\  \mathbf{x}^i_d.\ \ \text{Answer:}\  \mathbf{y}^i\ ”
\end{align}

\vspace*{-0.12in}
\paragraph{Reasoning instruction.} Our objective is to strengthen the reasoning capability of the LLM by explicitly incorporating reasoning instructions (refer to the blue annotations in Fig.~\ref{fig:prompt_for_rule}). These instructions commence with a statement reminiscent of the chain-of-thought methodology~\cite{wei2022chain}. The reasoning process is structured into two distinct phases. Initially, the LLM is prompted to hypothesize the causal connections or underlying patterns between the features and the overarching task, leveraging its general knowledge or common sense without consulting example demonstrations. This approach enables the LLM to utilize its prior understanding to contextualize the problem. Subsequently, in the second phase, the LLM synthesizes insights from example demonstrations along with the preliminary causal reasoning to construct specific rules associated with each class. This bifurcated reasoning paradigm helps the model avoid reliance on misleading associations found in extraneous columns and directs attention toward the most relevant features. To mitigate risks such as underfitting or unwieldy prompt sizes stemming from an excessive or insufficient number of rules, we constrain rule extraction for each class to a predefined count (i.e., 10)\footnote{The impact of rule quantity is discussed in Section~\ref{sec:analysis}.}. Furthermore, when creating rules, the LLM is instructed to consult the feature type information provided in the basic description to preclude any mismatches in rule formulation.

\vspace*{-0.12in}
\paragraph{Response instruction.} For more efficient parsing and subsequent utilization of the inferred rules, we prescribe a specific format for the LLM's output via detailed instructions (see yellow annotations in Fig.~\ref{fig:prompt_for_rule}). The prompt delineates both the expected response format for each answer class and the template to be used for rule formulation. By following these directives, the LLM is empowered to select a rule format corresponding to the nature of each feature. Absent further instruction, the LLM is capable of combining several rules using logical connectors such as AND or OR. A sample of the structured output can be found in Appendix~\ref{sec:outcome-rule}.

\subsection{Inferring Class Likelihood via Rules}
\label{Sec:training}

\paragraph{Rule-based feature creation via parsing.} In this stage, the rules identified earlier are employed to generate novel binary features. For every class, these features encode whether a sample fulfills the relevant rules for that class, assigning either ‘0’ or ‘1’. These features can subsequently inform predictions regarding a sample’s class membership. Nevertheless, as the LLM produces rules in natural language format, automated feature extraction necessitates a robust parsing approach. Despite efforts to standardize response and rule formatting during their generation to facilitate easier parsing, LLM outputs often exhibit unpredictable syntactic variation~\cite{kovavc2023large}, such as replacing symbols like ``$>$” or ``$<$” with their word equivalents ``greater than" or ``less than," which complicates the parsing process.

Rather than implementing intricate parsing scripts, we opt to use the LLM itself to handle noisy textual input. The LLM excels at grasping the semantic content of text in spite of syntactic alterations and is capable of generating concise code solutions with guidance, thanks to its exposure to code-based training data. To capitalize on these strengths, we format prompts with the relevant function names, input-output specifications, and the rules to be interpreted, and submit them to the LLM. Illustrative prompts and resulting outputs are presented in Figures~\ref{fig:prompt_for_parsing} and~\ref{sec:example_parsing}-\ref{sec:example_parse_outcome} in the Appendix. The resultant code is executed through Python’s \texttt{exec()} function for the required data transformation.

\vspace*{-0.12in}
\paragraph{Estimating the likelihood of class membership.} After acquiring binary features derived from class-specific rules, one straightforward approach to estimate the probability of a sample belonging to a class is by tallying how many of its binary features (per class) are satisfied—namely, summing the ‘1’ values associated with each class. Still, because certain rules have greater predictive value than others, it is important to determine their respective weights using training data. This can be achieved by fitting a standard ML model on each class’s set of binary features to learn their impact. For instance, one might employ a linear model with no intercept term. Specifically, the binary feature vector derived from sample $\mathbf{x}^i$ for class $k$ is denoted as $\mathbf{z}_k^i \in \{0, 1\}^R$, where $R$ represents the count of class-specific rules and $c$ is the total number of classes. The probability vector $\mathbf{p}^i$ for each class is subsequently calculated as:

\begin{align}
\text{logit}_k^i &= \max(\mathbf{w}_k, 0) \cdot \mathbf{z}_k^i, \nonumber \\
\mathbf{p}^i &= \text{Softmax}({[\text{logit}_{1}^i, ..., \text{logit}_{c}^i]}). \label{eq:infer_prob}
\end{align}

In this context, $\mathbf{w}_k$ denotes the tunable weights in the linear model, serving as indicators of the relevance of each input feature. By constraining these weights within a positive domain, we maintain that features contributed by the LLM cannot negatively influence the likelihood of a class during the training process. Subsequently, the logit outputs are transformed into class probabilities by applying the softmax operation. We employ cross-entropy loss minimization to optimize the weights $\mathbf{w}_k$, using a limited set of training examples in a few-shot learning scenario.

\vspace*{-0.12in}
\paragraph{Ensembling with bagging.} To conclude, we iteratively carry out the entire modeling procedure multiple times, producing a set of inference models whose predictions are aggregated in an ensemble manner. Considering the ensemble of trained weights across $T$ independent runs, expressed as $\{\mathbf{w}[t]\}_{t=1}^T$, final predictions are generated by averaging the class probabilities $\mathbf{p}^i$ obtained for each weight $\mathbf{w}[t]$, following Eq.~\ref{eq:infer_prob}.

To foster diversity among the rules generated for ensembling, we configure the LLM with an appropriately elevated temperature during inference. Additionally, we randomize the sequence of few-shot examples, based on findings that the order can significantly influence LLM response~\cite{lu2022fantastically}\footnote{Further investigation into rule diversity is available in Appendix~\ref{sec:diversity}}. Leveraging this prediction heterogeneity to strengthen the ensemble, we use bagging, which selects subsets of features or samples in each run; this strategy is particularly advantageous when facing large-scale training datasets or numerous features. This ensemble methodology yields two principal benefits. First, in scenarios where some LLM-generated rules are erroneous, others that are correct can mitigate these errors, thereby improving system resilience. This aligns with the principle of self-consistency in LLMs~\cite{wang2022self}, as rules recurrently produced across trials are statistically more trustworthy. Second, the approach alleviates prompt-length restrictions inherent to LLMs. When input tabular data exceeds permissible prompt sizes, bagging can limit input scope, promoting robust model performance. Though individual rule sets may fall short in covering the full range of feature dependencies, the composite ensemble of weak learners constructed over several trials diminishes the negative effects stemming from limited feature or sample representation within singular models.

\section{Experiments}
To assess the effectiveness of \textsf{FeatLLM}, we perform a series of evaluations using diverse tabular datasets, accompanied by several analyses to deepen understanding of the proposed method. Additional experimental results—including those with alternative LLMs, qualitative explorations, and other facets—are available in the Appendix due to page limitations.

\subsection{Performance Evaluation}
\paragraph{Datasets.}
Our empirical study spans 13 datasets, each targeting binary or multi-class classification objectives: (1) Adult~\cite{asuncion2007uci}, where the task is to determine if a person earns in excess of \$50,000 annually; (2) Bank~\cite{moro2014data}, focused on forecasting a client’s likelihood of subscribing to a term deposit; (3) Blood~\cite{yeh2009knowledge}, aiming to predict whether a blood donor will donate again; (4) Car~\cite{kadra2021well}, for classifying vehicle quality; (5) Communities~\cite{misc_communities_and_crime_183}, which estimates crime rates within particular areas; (6) Credit-g~\cite{kadra2021well}, used to evaluate whether someone represents a favorable or unfavorable credit risk; (7) Diabetes\footnote{\texttt{kaggle.com/datasets/uciml/pima-indians-diabetes-database}}, predicting an individual’s likelihood of having diabetes; (8) Heart\footnote{\texttt{kaggle.com/datasets/fedesoriano/heart-failure-prediction}}, which gauges the probability of coronary artery disease; and (9) Myocardial~\cite{misc_myocardial_infarction_complications_579}, designed to identify if a person is afflicted by chronic heart failure.

The dataset pool also includes recent additions and artificial constructs which are rarely utilized and did not contribute to the LLM pre-training corpus: (10) Cultivars, for forecasting soybean cultivar grain yield\footnote{\texttt{archive.ics.uci.edu/dataset/913}}; (11) NHANES, which categorizes individuals by age group based on health records\footnote{\texttt{archive.ics.uci.edu/dataset/887}}; (12) Sequence-type, concerned with categorizing sequences as arithmetic, geometric, Fibonacci, or Collatz; (13) Solution-mix, which predicts whether a mixture, formed from four solutions, contains a concentration exceeding 0.5. The Cultivars and NHANES datasets were submitted to the UCI Machine Learning Repository post-September 2023, guaranteeing exclusion from the training data for the LLM API (gpt-3.5-turbo-0613) employed in our methodology. The Sequence-type and Solution-mix datasets are synthetic, thus preventing potential memorization by the LLM API and allowing unbiased evaluation.

These datasets represent a wide spectrum in terms of scale and intricacy, as summarized in Table~\ref{Tab:basic-desc}.
Each dataset furnishes explicit names and explanations for all attributes. For example, datasets like `Communities' and `Myocardial' feature more than 100 attributes, introducing significant challenges, particularly given the finite context window length of LLMs.

\vspace*{-0.12in}
\paragraph{Baselines.}
For performance comparison, we benchmark against nine competing methods. The first set encompasses classical approaches for tabular prediction: (1) Logistic regression (LogReg), (2) XGBoost~\cite{chen2016xgboost}, and (3) RandomForest~\cite{ho1995random}. The subsequent pair presumes access to datasets lacking labels: (4) SCARF~\cite{bahri2022scarf} and (5) STUNT~\cite{nam2023stunt}. It is important to highlight the practical constraint in acquiring large collections of unlabeled data for many real-world cases. Another contemporary reference is (6) TabPFN~\cite{hollmann2022tabpfn}, which utilizes synthetically generated datasets with varied distributions for model pretraining. The final cohort is comprised of methods integrating LLMs: (7) In-context learning (In-context)~\cite{wei2022emergent}, (8) TABLET~\cite{slack2023tablet}, and (9) TabLLM~\cite{hegselmann2023tabllm}. In-context learning incorporates few-shot exemplars directly within prompts, foregoing parameter updates. TABLET augments the inference process by embedding supplemental cues via an external classifier that applies rule sets and prototypes. TabLLM leverages efficient tuning strategies such as IA3~\cite{liu2022few} to optimize the LLM with limited training samples.

\vspace*{-0.12in}
\paragraph{Implementation details.}
\textsf{FeatLLM} adopts GPT-3.5 as the underlying LLM, though it remains adaptable to alternative LLM configurations (refer to Appendix~\ref{sec:backbones} for comparative results with various backbones). The inference process uses a temperature of 0.5, while the top-p parameter retains its API default value of 1. For ensemble and rule extraction, we specify 20 ensembles and 10 rules. Further analysis regarding hyper-parameter selection can be found in Figure~\ref{fig:hyperparameter2} as well as Appendix~\ref{sec:hyper-parameter}. The linear model leverages the Adam optimizer, set at a learning rate of 0.01, and is trained over 200 epochs. Epoch selection is optimized through $k$-fold cross-validation. When reproducing baseline models, in-context learning is conducted using GPT-3.5; TabLLM follows its original protocol by utilizing T0~\cite{sanh2022multitask}. It is important to note that TabLLM limits LLM usage to models with publicly released checkpoints, hence GPT-3.5 is excluded from these experiments. For classic machine learning algorithms including LogReg, XGBoost, and RandomForest, parameters are fine-tuned by applying Grid-search in combination with $k$-fold cross-validation. The value for $k$ is chosen as 2 or 4 to ensure each class is represented at least once in the training data. Additional details on the implementation of baselines are provided in Appendix~\ref{sec:baselines}.

\vspace*{-0.12in}
\paragraph{Results.}
Tables~\ref{tab:main1} and \ref{tab:main2} display comparative performance results across various datasets. Each experiment was conducted three times, each with different splits of the training and test sets, and we report both the mean and standard deviation of the resulting AUC scores. Across these evaluations, \textsf{FeatLLM} reliably places first or second among baseline methods. Notably, our method achieves similar outcomes to traditional tabular learning models using just 4-shot samples, whereas those baselines require 32-shot samples for comparable results (refer to Fig.~\ref{fig:compares}). Although methods like STUNT and TabPFN recorded notable gains on some datasets, their overall performance increase relative to standard machine learning techniques was limited. Furthermore, LLM-based techniques, including in-context learning, TABLET, and TabLLM, encountered inference issues on datasets characterized by a high number of features. By integrating bagging and ensemble strategies, our approach remains effective and maintains strong results on such high-dimensional data. It should also be emphasized that the bagging and ensemble ideas presented here are compatible with other LLM-based methods; however, implementing them would require a doubling of per-sample inferences, considerably raising computational demands and rendering such modifications impractical.

\subsection{Analysis \& Discussion}
\label{sec:analysis}
\paragraph{Ablation study.} To understand the influence of the principal components in our model, we perform a series of ablation experiments by modifying: (1) \textsf{-Tuning}: removing the weight adjustment step for the linear model and using only a sum over binary features for logit computation; (2) \textsf{-Ensemble}: excluding the ensemble procedure; (3) \textsf{-Description}: leaving out feature descriptions; and (4) \textsf{-Reasoning}: skipping Step-1 in the rule-generation reasoning instructions.

Table~\ref{tab:ablation} provides a summary of the changes in average AUC values over all datasets resulting from each ablation. Performance declines for every case upon modifying an ablated element. An intriguing finding is that the value of weight tuning grows as the shot count increases, implying that larger sample sizes allow for more precise rule importance estimation, thereby heightening the impact of the tuning procedure. By contrast, leveraging feature descriptions and reasoning guidance becomes most advantageous when only a small number of examples is available. This underscores the role of exploiting prior LLM knowledge in scenarios with limited data. Finally, our ensemble strategy persistently enhances results across all tested configurations.

\vspace*{-0.12in}
\paragraph{Training \& inference time comparison.} We evaluate and contrast the training and inference durations across several methods. All experiments utilize the Adult dataset with a 16-shot training configuration. Unless otherwise stated, a single A100 GPU is used for running these experiments, with TabLLM requiring four A100 GPUs due to its model parallelism architecture. Table~\ref{tab:cost} outlines the measured runtimes. It is evident that \textsf{FeatLLM} achieves notably efficient inference times, matching those of traditional approaches. By comparison, LLM-based models like In-context learning, TABLET, and TabLLM incur considerably greater inference latency. With respect to training time, \textsf{FeatLLM} aligns with other LLM-centric techniques; specifically, it involves only 30 API queries, which are not financially demanding and can even be executed in parallel.

\vspace*{-0.12in}
\paragraph{Quality of features generated by \textsf{FeatLLM}.}
To assess the practical value of rule-based features produced by \textsf{FeatLLM}, we conduct an experiment focused on their informativeness for real-world tasks. Our evaluation involves calculating the average AUROC between each generated feature and the target labels, followed by determining the proportion of features whose AUROC surpasses 0.5 in each dataset. The results, presented in Table~\ref{tab:quality}, demonstrate that the majority of rules generated are strongly associated with the target attribute, providing valuable information for tackling the prediction task (see the ``Before tuning'' column). Furthermore, during the tuning phase with a linear model, features that exhibit low relevance to the target data (those whose learned weights fall below a set threshold) are automatically excluded from consideration (see the ``After tuning'' column). The threshold is fixed at 0.1, based on the observed distribution of feature weights.

\vspace*{-0.12in}
\paragraph{Effect of adding spurious correlations.} To evaluate the ability of different approaches to mitigate the influence of irrelevant or spurious correlations under low-shot conditions, we performed a dedicated analysis. Our study involved the Heart dataset, to which we appended randomly selected columns from the Adult dataset that have no relevance to the Heart task. We examined how model performance responded as the number of extraneous columns was incrementally increased. To ensure fairness, each method was provided with 8 labeled samples and no access to unlabeled data, which means techniques such as SCARF and STUNT were omitted from this evaluation.

Figure~\ref{fig:spurious} depicts how spurious correlations affect model accuracy. Among all compared methods, \textsf{FeatLLM} experienced the smallest drop in performance due to spurious features. This advantage likely stems from the framework’s requirement to explicitly align rules with the intended task using prior information, thereby avoiding reliance on irrelevant columns in rule formulation and resulting in increased robustness. Notably, less than 2\% of the extracted rules referenced unrelated columns. However, we also found that \textsf{FeatLLM} can occasionally internalize spurious correlations and misinterpret the causal relationships between task and features if columns from datasets with related content are randomly incorporated (see Appendix~\ref{sec:spurious-appendix}).

\vspace*{-0.12in}
\paragraph{Hyper-parameter analysis}
\textsf{FeatLLM} employs two primary hyper-parameters: the ensemble count and the quantity of rules per class generated from each LLM query. To assess their effects on performance, we performed a series of experiments. Findings indicate that increasing the ensemble number enhances performance up to a threshold (approximately 20 ensembles), after which improvements plateau (refer to Fig.~\ref{fig:appendix-hyperparameter1} in Appendix). Concerning rule count, we did not observe a statistically significant variance in outcomes, although selecting 10 rules struck an optimal balance between prompt length and accuracy (see Fig.~\ref{fig:hyperparameter2}). We conjecture that increasing rule quantity augments the input’s dimensionality—supplying additional information—but also exacerbates overfitting, especially in settings with limited data.

\section{Conclusion}
We introduce \textsf{FeatLLM}, an advancement in LLM-based few-shot tabular learning techniques. Rather than relying solely on LLMs for sample-level predictions, \textsf{FeatLLM} emphasizes synthesizing decision criteria by adopting feature engineering methodologies. The architecture leverages rule generation and a bagging-based ensemble, requiring minimal inference overhead and operating exclusively with API calls—no model training required—thereby circumventing restrictions associated with data feature scale. This approach results in robust predictive performance and serves as an effective, data-efficient solution for applying LLMs to real-world tabular datasets.

\textsf{FeatLLM} is purposely tailored for tasks characterized by scant labeled data. In future work, we aim to extend its feature creation strategies to settings with more abundant samples, incorporate various feature constructs beyond rules, and enhance its interpretability, allowing it to be utilized in broader contexts.

\section{Broader Impact}
The methodological framework proposed, \textsf{FeatLLM}, leverages the existing domain knowledge in LLMs to elevate performance in tabular learning beyond conventional supervised paradigms. Our research addresses challenges related to overfitting arising from small training sets by integrating LLM-driven prior knowledge to enhance data efficiency. This is especially advantageous in domains such as finance and healthcare, where label acquisition is expensive or difficult and where pretrained LLMs encompass pertinent textual knowledge. As a prospective avenue for broader implementation, it becomes crucial to scrutinize how societal biases or inaccuracies contained within LLM priors may influence predictions—potentially overtaking empirical data—and this warrants systematic analysis moving forward.

\end{document}